(1) 正入射两非磁性介质的电场反射系数为 $(n_i-n_j)/(n_i+n_j)$。本题 $n_0=n_2=1$，于是 $r_{01}=(1-n)/(1+n)\equiv r_0$、$r_{12}=(n-1)/(n+1)=-r_0$；两个反射面具有相反的相位翻转性质。
(2) 每次薄膜内往返积累相位 $2\delta$，结合界面多重反射可得 $r=(r_{01}+r_{12}e^{2i\delta})/(1+r_{01}r_{12}e^{2i\delta})$。代入两侧均为空气时 $r=r_0(1-e^{2i\delta})/(1-r_0^2e^{2i\delta})$，该表达式保留所有反射次数和干涉相位。
(3) $e^{2i\delta}=1$ 时分子为零，即 $\delta=m\pi$，故 $2nd=m\lambda$，其中非负整数 $m$ 对应半波厚度。若 $\delta=(m+\frac12)\pi$，则 $r=2r_0/(1+r_0^2)$、$R=4r_0^2/(1+r_0^2)^2$；$n\to1$ 时 $r_0\to0$，整膜在任何厚度都无反射。
(4) 薄膜两侧介质都是空气，入射与出射行波具有相同光学导纳。介质无吸收并且稳态无净能量储存，故 $R+T=1$，即 $T=1-|r|^2$。若仅把前后界面反射率直接相加，则丢失交叉干涉项及无限次反射，甚至可能违反精确能流守恒。
